\documentclass[a4paper,11pt]{article}
\usepackage{ctex}
\usepackage{amsmath, amssymb}
\usepackage{geometry}
\geometry{left=2.5cm, right=2.5cm, top=2.5cm, bottom=2.5cm}

\begin{document}

\begin{center}
    \Large\textbf{上机实验报告：Gauss 消去法及其数值稳定性分析}
\end{center}

\section{实验问题与理论精确解}
针对 84 阶三对角线性方程组 $\mathbf{A}x = b$，其系数矩阵主对角线元素为 6，次下对角线为 8，次上对角线为 1。右端项向量首位为 7，末位为 14，其余各项均为 15。

通过观察原方程组系数，将全 1 向量 $x^* = [1, 1, \dots, 1]^T$ 代入方程组进行检验：
\begin{itemize}
    \item 第 1 行：$6 \times 1 + 1 \times 1 = 7 = b_1$
    \item 第 $i$ 行 ($1 < i < 84$)：$8 \times 1 + 6 \times 1 + 1 \times 1 = 15 = b_i$
    \item 第 84 行：$8 \times 1 + 6 \times 1 = 14 = b_{84}$
\end{itemize}
上述条件均严格成立，故该方程组的理论精确解为全 1 向量。

\section{数值实验结果对比}
采用不选主元和列主元的 Gauss 消去法求解该方程组，提取后 10 项计算结果并与理论精确解进行无穷范数误差对比，结果如下：

\subsection{不选主元 Gauss 消去法}
\begin{itemize}
    \item 数值解（后 10 项）：$[-2.09 \times 10^6, \, 4.18 \times 10^6, \, -8.35 \times 10^6, \, 1.66 \times 10^7, \, -3.30 \times 10^7, \, 6.50 \times 10^7, \, -1.25 \times 10^8, \, 2.34 \times 10^8, \, -4.02 \times 10^8, \, 5.36 \times 10^8]$
    \item 最大绝对误差：$5.368381 \times 10^8$
\end{itemize}

\subsection{列主元 Gauss 消去法}
\begin{itemize}
    \item 数值解（后 10 项）：$[1.0, \, 1.0, \, 1.0, \, 1.0, \, 1.0, \, 1.0, \, 1.000001, \, 0.999999, \, 1.000002, \, 0.999997]$
    \item 最大绝对误差：$2.797445 \times 10^{-6}$
\end{itemize}

\section{结果分析与对 Gauss 消去法的看法}
实验结果表明，对该特定方程组使用原生 Gauss 消去法会引发严重的数值崩溃，而列主元法则表现出极高的数值稳定性。具体分析如下：

\subsection{不选主元法的误差放大机制}
在不选主元的消元过程中，主对角元初始为 6，下方元素为 8，第一步消元乘子 $m_1 = 8/6 \approx 1.33$。随着消元的进行，主对角元更新公式为：
\begin{equation}
a_{kk} = 6 - \frac{8}{a_{k-1, k-1}}
\end{equation}
计算该序列的极限 $\lim_{k \to \infty} a_{kk} = L$，有 $L = 6 - 8/L$，解得稳态值 $L = 4$。

这意味着对角元会迅速递减并趋近于 4，而对应的消元乘子 $m_k = 8/a_{kk}$ 会逐渐增大并趋近于 2。由于乘子 $|m_k| > 1$，计算机浮点数表示下产生的微小舍入误差在每次迭代中被不断放大。经过 84 次消元后，误差呈指数级爆炸，最终使得最大误差达到了 $10^8$ 数量级，计算结果完全失真。

\subsection{选主元策略的数值稳定性}
引入列主元策略后，算法在每步消元前会强制选取当前列对角线及其下方绝对值最大的元素（即 8）作为主元。通过行交换，对角元变为 8，消元乘子变为 $m_k = 6/8 = 0.75$。

由于满足 $|m_k| < 1$，舍入误差在传递过程中被不断缩小和抑制，避免了指数级累积。实验中列主元法的最大误差稳定在 $2.79 \times 10^{-6}$ 量级，在当前浮点精度允许的范围内高度还原了精确解。

\subsection{结论}
原生的 Gauss 消去法在有限字长的计算机浮点运算环境中，本质上是一种数值不稳定的算法。在实际求解一般的大型线性方程组时，除非系数矩阵具备严格对角占优或对称正定等特殊性质，否则必须采用列主元或全主元策略。选主元的核心目的在于将消元乘子的绝对值控制在 1 以内，从而有效抑制舍入误差的传播，确保数值计算结果的可靠性。

\end{document}